\begin{ouverture}[college3]{13}{Circuits : intensités et tensions}{Partie 3 · L’énergie, ses transferts et ses conversions}{Quelles lois permettent de contrôler un montage ?}
\objectif{Schématiser série et dérivation.}
\objectif{Brancher les appareils et convertir les mesures.}
\objectif{Exploiter les lois des intensités et des tensions.}
\prerequis{Symboles usuels et circuit fermé.}
\prerequis{A, mA et V.}
\end{ouverture}
\pcPage{Le cours}{Circuits : intensités et tensions}
\phantomsection\label{ess-13}
\essSection{1}{Lire la topologie}
Deux dipôles sont en série s’ils appartiennent à une même branche sans dérivation entre eux. Deux dipôles sont en dérivation s’ils sont reliés aux mêmes deux nœuds. Le schéma indique les connexions ; la position sur la feuille ne suffit pas. Une source alimente le circuit fermé. Un court-circuit peut provoquer une intensité élevée et un échauffement : on coupe l’alimentation avant toute modification.\par
\pcFigure{figures/b13-circuits.pdf}{Topologie de deux circuits : suivre les connexions et repérer les deux nœuds.}
\essSection{2}{Mesurer correctement}
Un ampèremètre mesure une intensité I et se branche en série dans la branche étudiée. Un voltmètre mesure une tension U et se branche en dérivation aux bornes du dipôle. Utiliser les bornes et le calibre adaptés, commencer par un calibre suffisant et respecter le sens de connexion en courant continu. $1\unite A=1\,000\unite{mA}$. Un affichage négatif peut venir d’une inversion des bornes.\par
\clearpage
\essSection{3}{Les lois des intensités}
Dans une unique branche de dipôles en série, l’intensité du courant est la même en tout point en régime permanent. À un nœud, la somme des intensités entrantes égale celle des intensités sortantes ; avec deux branches réceptrices, $I_G=I_1+I_2$. Le courant n’est pas « consommé » par une lampe : c’est l’énergie qui lui est transférée.\par
\pcFigure{figures/b13-mesures.pdf}{Ampèremètre en série et voltmètre en dérivation : connexions distinctes.}
\essSection{4}{Les lois des tensions}
Dans une boucle simple comprenant la source et deux récepteurs en série, avec les valeurs positives et conventions indiquées, $U_G=U_1+U_2$. Deux dipôles reliés aux mêmes deux nœuds ont la même tension : $U_1=U_2=U_G$ dans le montage en dérivation directement alimenté. Ne pas confondre égalité des tensions et égalité des intensités : les courants des branches peuvent différer.\par
\begin{pcEncadre}{method}{Un exemple expliqué}Deux lampes en dérivation sont soumises à 6,0 V. Les intensités sont 120 mA et 180 mA. La source débite $I_G=0{,}120+0{,}180=0{,}300\unite A$. Chaque lampe reçoit 6,0 V ; cette égalité n’oblige pas les deux courants à être égaux.\end{pcEncadre}
\pcPage{Expérimenter et retenir}{Circuits : intensités et tensions}
\phantomsection\label{activite-13}
\begin{pcEncadre}{method}{Investigation en classe}Réaliser sous contrôle un circuit scolaire à deux lampes. Mesurer les intensités en série, puis les trois intensités d’un montage en dérivation. Contrôler ensuite les tensions sur les mêmes nœuds. Comparer aux lois, en tenant compte du calibre, de la lecture et de la stabilité de la source.\end{pcEncadre}
\pcSousTitre{Ce que je dois savoir expliquer}\begin{questions}
\item Choisir le bon montage de mesure.
\item Écrire la loi avant de remplacer les valeurs.
\item Vérifier la cohérence d’une mesure avec les connexions.
\end{questions}
\begin{pcEncadre}{attention}{Une erreur à éviter}Brancher un ampèremètre directement aux bornes d’une source peut créer un court-circuit. Il ne remplace jamais un voltmètre.\end{pcEncadre}
\pcRepere{cor-activite-13}{Pistes pour l’investigation}\pcGrille
\clearpage\markright{Je démarre}
\rubrique{Je démarre}
\par\noindent\begin{minipage}{\linewidth}
\begin{exercice}[1]{Le branchement}
\begin{questions}
\item Comment brancher un ampèremètre ?
\item Comment brancher un voltmètre ?
\end{questions}
\end{exercice}
\end{minipage}
\par\noindent\begin{minipage}{\linewidth}
\begin{exercice}[1]{Conversions}
\begin{questions}
\item Convertir 250 mA en A.
\item Convertir 0,032 A en mA.
\end{questions}
\end{exercice}
\end{minipage}
\par\noindent\begin{minipage}{\linewidth}
\begin{exercice}[1]{En série}
\begin{questions}
\item Dans une seule branche, on mesure 0,18 A avant une lampe. Quelle intensité après ?
\item Le courant est-il consommé ?
\end{questions}
\end{exercice}
\end{minipage}
\par\noindent\begin{minipage}{\linewidth}
\begin{exercice}[1]{Deux nœuds}
\pcFigureEx{figures/ex-circuit-parallele.pdf}{}
\begin{questions}
\item Deux lampes sont reliées aux mêmes deux nœuds. Comment nommer l’association ?
\item Quelle grandeur a la même valeur aux bornes des lampes ?
\end{questions}
\end{exercice}
\end{minipage}
\par\noindent\begin{minipage}{\linewidth}
\begin{exercice}[1]{Une somme}
\pcFigureEx{figures/rev-ex-noeud-035.pdf}{}
\begin{questions}
\item Deux courants de 0,10 A et 0,25 A sortent d’un nœud. Quelle intensité totale entre ?
\item Donner la valeur en mA.
\end{questions}
\end{exercice}
\end{minipage}
\clearpage\markright{Je m’entraîne}
\rubrique{Je m’entraîne}
\par\noindent\begin{minipage}{\linewidth}
\begin{exercice}[2]{Une branche inconnue}
\pcFigureEx{figures/rev-ex-noeud-080.pdf}{}
\begin{questions}
\item $I_G=0{,}80\unite A$, $I_1=250\unite{mA}$. Calculer I2.
\item Quelle erreur ferait-on en calculant $0{,}80-250$ ?
\end{questions}
\end{exercice}
\end{minipage}
\par\noindent\begin{minipage}{\linewidth}
\begin{exercice}[2]{Tensions en série}
\begin{questions}
\item Une source de 9 V alimente deux récepteurs en série. U1 = 3,5 V. Trouver U2.
\item Où brancher le voltmètre pour mesurer U2 ?
\end{questions}
\end{exercice}
\end{minipage}
\par\noindent\begin{minipage}{\linewidth}
\begin{exercice}[2]{Tensions en dérivation}
\begin{questions}
\item Deux résistances sont en dérivation sur une source idéale de 12 V. Donner leurs tensions.
\item Leurs intensités sont-elles forcément égales ?
\end{questions}
\end{exercice}
\end{minipage}
\par\noindent\begin{minipage}{\linewidth}
\begin{exercice}[2]{Une inversion}
\begin{questions}
\item Un voltmètre affiche −4,5 V. Quelle cause de connexion envisager ?
\item L’intensité est-elle pour autant forcément nulle ?
\end{questions}
\end{exercice}
\end{minipage}
\par\noindent\begin{minipage}{\linewidth}
\begin{exercice}[2]{Commander séparément}
\begin{questions}
\item Deux lampes doivent être commandées indépendamment. Quelle association proposer ?
\item Où placer les deux interrupteurs ?
\end{questions}
\end{exercice}
\end{minipage}
\clearpage\markright{Je raisonne}
\rubrique{Je raisonne}
\par\noindent\begin{minipage}{\linewidth}
\begin{exercice}[3]{Une mesure incohérente}
\begin{questions}
\item En série, un élève mesure 0,20 A puis 0,12 A. La loi est-elle invalidée par ces deux nombres seuls ?
\item Citer deux contrôles.
\end{questions}
\end{exercice}
\end{minipage}
\par\noindent\begin{minipage}{\linewidth}
\begin{exercice}[3]{Un nœud à trois sorties}
\pcFigureEx{figures/rev-ex-noeud-090.pdf}{}
\begin{questions}
\item Entrée : 0,90 A ; sorties connues : 0,20 A et 0,35 A. Calculer la troisième.
\item La somme des sorties doit-elle être plus petite que l’entrée en régime permanent ?
\end{questions}
\end{exercice}
\end{minipage}
\par\noindent\begin{minipage}{\linewidth}
\begin{exercice}[3]{Une lampe retirée}
\begin{questions}
\item Dans un montage en dérivation, on ouvre une branche sans toucher l’autre. L’autre lampe peut-elle rester allumée ?
\item Quelle hypothèse sur la source faut-il annoncer pour conserver exactement sa tension ?
\end{questions}
\end{exercice}
\end{minipage}
\clearpage
\par\noindent\begin{minipage}{\linewidth}
\begin{exercice}[3]{Court-circuit}
\pcFigureEx{figures/ex-lampe-court-circuitee.pdf}{Défaut à analyser sur le schéma, sans le réaliser.}
\begin{questions}
\item Un fil relie directement les deux bornes d’une lampe. Pourquoi la situation est-elle différente d’une branche avec une lampe ?
\item Que faire avant de modifier le montage ?
\end{questions}
\end{exercice}
\end{minipage}
\par\noindent\begin{minipage}{\linewidth}
\begin{exercice}[3]{Un cahier des charges}
\begin{questions}
\item Deux lampes doivent recevoir chacune 6 V, l’une sous 0,10 A et l’autre sous 0,20 A. Quelle association et quelle intensité de source prévoir ?
\item Pourquoi une association en série sur 6 V ne répond-elle pas en général à la demande ?
\end{questions}
\end{exercice}
\end{minipage}
\pcPage{Défi du chapitre}{Circuits : intensités et tensions}
\phantomsection\label{defi-ex-13}
\pcDefiTitre{13}{Le montage à vérifier}
\pcFigureEx{figures/rev-ex-montage-9V.pdf}{}
\begin{questions}
\item Deux branches sont reliées à une source de 9 V. La première comprend deux lampes en série sous 4 V et 5 V ; la seconde une résistance. Vérifier les tensions.
\item Les courants de branches valent 150 mA et 60 mA. Trouver IG.
\item Décrire un protocole de mesure de la tension de la résistance et du courant de la première branche.
\end{questions}
\pcGrille
